\documentclass[border=10pt,tikz]{standalone}
\usepackage{tikz}
\usepackage{kotex}
\setmainhangulfont{Apple SD Gothic Neo}[BoldFont={Apple SD Gothic Neo Bold}]
\setsanshangulfont{Apple SD Gothic Neo}[BoldFont={Apple SD Gothic Neo Bold}]
\usetikzlibrary{positioning,arrows.meta,shapes,fit,calc,backgrounds,decorations.pathreplacing}
\begin{document}
% _design.tex — shared TikZ design system for all blog diagrams
% Inputs nothing on its own — meant to be \input{} from each diagram .tex
% AFTER \begin{document}, BEFORE \begin{tikzpicture}.
%
% Companion files:
%   _design-math.tex — math diagrams (PGFPlots, book-notes palette)
%
% Provides a unified visual language across GoF / DSA / EMC++ / Beautiful / 시리즈:
%   - Colors: muted, accessible, color-coded by role
%   - Node styles: consistent sizing, rounded, subtle borders
%   - Edge styles: UML-style inheritance, aggregation, plus dashed alternates
%   - Annotation: callouts, notes, badges
%
% Usage:
%   \documentclass[border=10pt,tikz]{standalone}
%   \usepackage{tikz}
%   \usetikzlibrary{positioning,arrows.meta,shapes,fit,backgrounds,calc,decorations.pathreplacing}
%   \begin{document}
%   % _design.tex — shared TikZ design system for all blog diagrams
% Inputs nothing on its own — meant to be \input{} from each diagram .tex
% AFTER \begin{document}, BEFORE \begin{tikzpicture}.
%
% Companion files:
%   _design-math.tex — math diagrams (PGFPlots, book-notes palette)
%
% Provides a unified visual language across GoF / DSA / EMC++ / Beautiful / 시리즈:
%   - Colors: muted, accessible, color-coded by role
%   - Node styles: consistent sizing, rounded, subtle borders
%   - Edge styles: UML-style inheritance, aggregation, plus dashed alternates
%   - Annotation: callouts, notes, badges
%
% Usage:
%   \documentclass[border=10pt,tikz]{standalone}
%   \usepackage{tikz}
%   \usetikzlibrary{positioning,arrows.meta,shapes,fit,backgrounds,calc,decorations.pathreplacing}
%   \begin{document}
%   % _design.tex — shared TikZ design system for all blog diagrams
% Inputs nothing on its own — meant to be \input{} from each diagram .tex
% AFTER \begin{document}, BEFORE \begin{tikzpicture}.
%
% Companion files:
%   _design-math.tex — math diagrams (PGFPlots, book-notes palette)
%
% Provides a unified visual language across GoF / DSA / EMC++ / Beautiful / 시리즈:
%   - Colors: muted, accessible, color-coded by role
%   - Node styles: consistent sizing, rounded, subtle borders
%   - Edge styles: UML-style inheritance, aggregation, plus dashed alternates
%   - Annotation: callouts, notes, badges
%
% Usage:
%   \documentclass[border=10pt,tikz]{standalone}
%   \usepackage{tikz}
%   \usetikzlibrary{positioning,arrows.meta,shapes,fit,backgrounds,calc,decorations.pathreplacing}
%   \begin{document}
%   \input{../_design.tex}     % adjust relative path
%   \begin{tikzpicture}[blog]  % apply 'blog' style
%   ...
%
% Backward-compat note: This file only ADDS to the public API. Existing styles
% (abs/con/imp/cli/hi/vx/q/inh/agg/uses/arr/edge/alt/ret/thr/group/glabel) are
% preserved. New helpers are documented below.

% ─── Core palette (fills paired with darker borders) ─────────
\definecolor{absfill}{RGB}{220,232,247}   % cool blue
\definecolor{absbord}{RGB}{72,118,182}
\definecolor{confill}{RGB}{249,222,222}   % warm red
\definecolor{conbord}{RGB}{197,93,93}
\definecolor{impfill}{RGB}{223,238,222}   % soft green
\definecolor{impbord}{RGB}{99,154,93}
\definecolor{clifill}{RGB}{251,239,205}   % pale yellow
\definecolor{clibord}{RGB}{200,167,72}
\definecolor{hifill} {RGB}{253,224,200}   % accent orange
\definecolor{hibord} {RGB}{210,135,68}
\definecolor{nodefill}{RGB}{250,250,250}
\definecolor{nodebord}{RGB}{120,120,120}
\definecolor{edgec}  {RGB}{80,80,80}
\definecolor{edgesc} {RGB}{145,145,145}
\definecolor{groupbord}{RGB}{200,200,200}

% ─── Extended palette (new) ──────────────────────────────────
\definecolor{prpfill}{RGB}{233,224,247}   % muted purple (state / interface)
\definecolor{prpbord}{RGB}{132,103,182}
\definecolor{tealfill}{RGB}{215,235,235}  % teal (services / utility)
\definecolor{tealbord}{RGB}{ 78,135,140}
\definecolor{warnfill}{RGB}{251,229,229}  % deeper warn (error states)
\definecolor{warnbord}{RGB}{180, 55, 55}
\definecolor{okfill}  {RGB}{220,240,225}  % success
\definecolor{okbord}  {RGB}{ 75,140, 90}
\definecolor{neutfill}{RGB}{238,238,238}  % neutral panel
\definecolor{neutbord}{RGB}{170,170,170}
\definecolor{inkc}    {RGB}{ 30, 30, 30}  % strong ink (titles)
\definecolor{mutec}   {RGB}{120,120,120}  % muted text
\definecolor{focusc}  {RGB}{210, 70, 70}  % focus / highlight accent
\definecolor{gridc}   {RGB}{220,220,220}  % background grid

\tikzset{
  blog/.style={
    font=\sffamily\small\linespread{1.18}\selectfont,
    every node/.style={font=\sffamily\small\linespread{1.18}\selectfont}
  },
  % ─── existing node roles (unchanged) ───
  cls/.style={
    rectangle, rounded corners=2pt,
    draw=nodebord, line width=0.4pt,
    minimum height=8mm, minimum width=28mm,
    inner sep=3pt, fill=nodefill, align=center
  },
  abs/.style={cls, draw=absbord, fill=absfill, font=\sffamily\itshape\small},
  con/.style={cls, draw=conbord, fill=confill},
  imp/.style={cls, draw=impbord, fill=impfill},
  cli/.style={cls, draw=clibord, fill=clifill},
  hi/.style={cls, draw=hibord, fill=hifill, font=\sffamily\bfseries\small},
  % ─── new node roles ───
  prp/.style={cls, draw=prpbord, fill=prpfill},               % protocol / interface / state
  tealn/.style={cls, draw=tealbord, fill=tealfill},           % service / utility
  warn/.style={cls, draw=warnbord, fill=warnfill, font=\sffamily\bfseries\small},  % error state
  ok/.style={cls, draw=okbord, fill=okfill},                  % success state
  neut/.style={cls, draw=neutbord, fill=neutfill},            % neutral
  % ─── size variants ───
  clsSm/.style={cls, minimum height=6mm, minimum width=20mm, font=\sffamily\footnotesize, inner sep=2pt},
  clsLg/.style={cls, minimum height=11mm, minimum width=36mm, font=\sffamily\normalsize},
  absSm/.style={clsSm, draw=absbord, fill=absfill, font=\sffamily\itshape\footnotesize},
  conSm/.style={clsSm, draw=conbord, fill=confill},
  impSm/.style={clsSm, draw=impbord, fill=impfill},
  % ─── circle (existing) ───
  vx/.style={
    circle, draw=absbord, line width=0.4pt,
    minimum size=8mm, fill=absfill, inner sep=0pt
  },
  vxc/.style={vx, draw=conbord, fill=confill},
  vxh/.style={vx, draw=hibord, fill=hifill, font=\sffamily\bfseries},
  % new circle variants
  vxi/.style={vx, draw=impbord, fill=impfill},                % impl green
  vxp/.style={vx, draw=prpbord, fill=prpfill},                % protocol purple
  vxw/.style={vx, draw=warnbord, fill=warnfill},              % warn red
  vxn/.style={vx, draw=nodebord, fill=nodefill},              % neutral
  vxSm/.style={vx, minimum size=6mm, font=\sffamily\footnotesize},
  vxLg/.style={vx, minimum size=11mm, font=\sffamily\normalsize},
  % ─── decision diamond (existing) ───
  q/.style={
    diamond, draw=clibord, line width=0.4pt, fill=clifill,
    inner sep=2pt, aspect=2.4, align=center
  },
  % rounded callout / note (new)
  note/.style={
    rectangle, rounded corners=3pt,
    draw=clibord, line width=0.4pt, fill=clifill,
    inner sep=3pt, font=\sffamily\footnotesize, align=left
  },
  badge/.style={                                              % small inline marker
    rectangle, rounded corners=2pt,
    draw=hibord, line width=0.3pt, fill=hifill,
    inner sep=1.5pt, font=\sffamily\scriptsize\bfseries
  },
  % ─── edges (existing) ───
  inh/.style={                                  % UML inheritance
    -{Triangle[length=3mm,width=3mm,open]},
    draw=edgec, line width=0.5pt
  },
  agg/.style={                                  % UML aggregation
    {Diamond[length=2.5mm,width=2.5mm,open]}-{Stealth[length=2.2mm]},
    draw=edgec, line width=0.45pt
  },
  uses/.style={                                 % directed reference
    -{Stealth[length=2.2mm]}, draw=edgec, line width=0.5pt
  },
  arr/.style={                                  % plain arrow (graphs/trees)
    -{Stealth[length=2mm]}, draw=edgec, line width=0.45pt
  },
  edge/.style={                                 % undirected (trees)
    draw=edgec, line width=0.45pt
  },
  alt/.style={                                  % dashed: similar/alt
    -{Stealth[length=2mm]}, dashed, draw=edgesc, line width=0.45pt
  },
  ret/.style={                                  % dashed creates/returns
    -{Stealth[length=2mm]}, dashed, draw=edgesc, line width=0.45pt
  },
  thr/.style={                                  % threading link (red dashed)
    ->, dashed, draw=conbord, line width=0.5pt
  },
  % ─── new edges ───
  comp/.style={                                 % UML composition (filled diamond)
    {Diamond[length=2.5mm,width=2.5mm]}-{Stealth[length=2.2mm]},
    draw=edgec, line width=0.45pt
  },
  realize/.style={                              % UML realize (impl interface)
    -{Triangle[length=3mm,width=3mm,open]}, dashed,
    draw=edgec, line width=0.5pt
  },
  msgsync/.style={                              % synchronous message (sequence)
    -{Stealth[length=2.6mm]}, draw=edgec, line width=0.55pt
  },
  msgasync/.style={                             % asynchronous message
    -{Stealth[length=2.4mm,inset=0.4mm,fill=edgec]},
    draw=edgec, line width=0.5pt
  },
  msgret/.style={                               % return message (dashed)
    -{Stealth[length=2.2mm]}, dashed, draw=edgesc, line width=0.45pt
  },
  flow/.style={                                 % control flow (thicker, hi-vis)
    -{Stealth[length=2.4mm]}, draw=inkc, line width=0.7pt
  },
  focus/.style={                                % focus / hot path
    -{Stealth[length=2.4mm]}, draw=focusc, line width=0.8pt
  },
  selfloop/.style={                             % self loop on a node
    ->, draw=edgec, line width=0.45pt, in=120, out=60, loop, looseness=8
  },
  % ─── group / region (existing + new) ───
  group/.style={
    rectangle, rounded corners=4pt,
    draw=groupbord, line width=0.4pt,
    inner sep=10pt, fill=black!2
  },
  glabel/.style={font=\sffamily\bfseries\small, anchor=north west, inner sep=2pt},
  panel/.style={                                % stronger panel for layered diagrams
    rectangle, rounded corners=4pt,
    draw=neutbord, line width=0.45pt,
    inner sep=10pt, fill=neutfill!50
  },
  panelTitle/.style={font=\sffamily\bfseries\small, text=inkc, anchor=north west, inner sep=2pt},
  panelSub/.style={font=\sffamily\footnotesize, text=mutec, anchor=north west, inner sep=2pt},
  % ─── typography ───
  caption/.style={font=\sffamily\footnotesize\itshape, text=mutec, align=center},
  legendkey/.style={rectangle, minimum width=10mm, minimum height=3mm, inner sep=0pt, line width=0.3pt},
  % ─── grid / coordinate hints ───
  bggrid/.style={draw=gridc, line width=0.3pt, dotted},
}

% Korean font convenience macro — included by every Hangul diagram
\providecommand{\KOR}{}
     % adjust relative path
%   \begin{tikzpicture}[blog]  % apply 'blog' style
%   ...
%
% Backward-compat note: This file only ADDS to the public API. Existing styles
% (abs/con/imp/cli/hi/vx/q/inh/agg/uses/arr/edge/alt/ret/thr/group/glabel) are
% preserved. New helpers are documented below.

% ─── Core palette (fills paired with darker borders) ─────────
\definecolor{absfill}{RGB}{220,232,247}   % cool blue
\definecolor{absbord}{RGB}{72,118,182}
\definecolor{confill}{RGB}{249,222,222}   % warm red
\definecolor{conbord}{RGB}{197,93,93}
\definecolor{impfill}{RGB}{223,238,222}   % soft green
\definecolor{impbord}{RGB}{99,154,93}
\definecolor{clifill}{RGB}{251,239,205}   % pale yellow
\definecolor{clibord}{RGB}{200,167,72}
\definecolor{hifill} {RGB}{253,224,200}   % accent orange
\definecolor{hibord} {RGB}{210,135,68}
\definecolor{nodefill}{RGB}{250,250,250}
\definecolor{nodebord}{RGB}{120,120,120}
\definecolor{edgec}  {RGB}{80,80,80}
\definecolor{edgesc} {RGB}{145,145,145}
\definecolor{groupbord}{RGB}{200,200,200}

% ─── Extended palette (new) ──────────────────────────────────
\definecolor{prpfill}{RGB}{233,224,247}   % muted purple (state / interface)
\definecolor{prpbord}{RGB}{132,103,182}
\definecolor{tealfill}{RGB}{215,235,235}  % teal (services / utility)
\definecolor{tealbord}{RGB}{ 78,135,140}
\definecolor{warnfill}{RGB}{251,229,229}  % deeper warn (error states)
\definecolor{warnbord}{RGB}{180, 55, 55}
\definecolor{okfill}  {RGB}{220,240,225}  % success
\definecolor{okbord}  {RGB}{ 75,140, 90}
\definecolor{neutfill}{RGB}{238,238,238}  % neutral panel
\definecolor{neutbord}{RGB}{170,170,170}
\definecolor{inkc}    {RGB}{ 30, 30, 30}  % strong ink (titles)
\definecolor{mutec}   {RGB}{120,120,120}  % muted text
\definecolor{focusc}  {RGB}{210, 70, 70}  % focus / highlight accent
\definecolor{gridc}   {RGB}{220,220,220}  % background grid

\tikzset{
  blog/.style={
    font=\sffamily\small\linespread{1.18}\selectfont,
    every node/.style={font=\sffamily\small\linespread{1.18}\selectfont}
  },
  % ─── existing node roles (unchanged) ───
  cls/.style={
    rectangle, rounded corners=2pt,
    draw=nodebord, line width=0.4pt,
    minimum height=8mm, minimum width=28mm,
    inner sep=3pt, fill=nodefill, align=center
  },
  abs/.style={cls, draw=absbord, fill=absfill, font=\sffamily\itshape\small},
  con/.style={cls, draw=conbord, fill=confill},
  imp/.style={cls, draw=impbord, fill=impfill},
  cli/.style={cls, draw=clibord, fill=clifill},
  hi/.style={cls, draw=hibord, fill=hifill, font=\sffamily\bfseries\small},
  % ─── new node roles ───
  prp/.style={cls, draw=prpbord, fill=prpfill},               % protocol / interface / state
  tealn/.style={cls, draw=tealbord, fill=tealfill},           % service / utility
  warn/.style={cls, draw=warnbord, fill=warnfill, font=\sffamily\bfseries\small},  % error state
  ok/.style={cls, draw=okbord, fill=okfill},                  % success state
  neut/.style={cls, draw=neutbord, fill=neutfill},            % neutral
  % ─── size variants ───
  clsSm/.style={cls, minimum height=6mm, minimum width=20mm, font=\sffamily\footnotesize, inner sep=2pt},
  clsLg/.style={cls, minimum height=11mm, minimum width=36mm, font=\sffamily\normalsize},
  absSm/.style={clsSm, draw=absbord, fill=absfill, font=\sffamily\itshape\footnotesize},
  conSm/.style={clsSm, draw=conbord, fill=confill},
  impSm/.style={clsSm, draw=impbord, fill=impfill},
  % ─── circle (existing) ───
  vx/.style={
    circle, draw=absbord, line width=0.4pt,
    minimum size=8mm, fill=absfill, inner sep=0pt
  },
  vxc/.style={vx, draw=conbord, fill=confill},
  vxh/.style={vx, draw=hibord, fill=hifill, font=\sffamily\bfseries},
  % new circle variants
  vxi/.style={vx, draw=impbord, fill=impfill},                % impl green
  vxp/.style={vx, draw=prpbord, fill=prpfill},                % protocol purple
  vxw/.style={vx, draw=warnbord, fill=warnfill},              % warn red
  vxn/.style={vx, draw=nodebord, fill=nodefill},              % neutral
  vxSm/.style={vx, minimum size=6mm, font=\sffamily\footnotesize},
  vxLg/.style={vx, minimum size=11mm, font=\sffamily\normalsize},
  % ─── decision diamond (existing) ───
  q/.style={
    diamond, draw=clibord, line width=0.4pt, fill=clifill,
    inner sep=2pt, aspect=2.4, align=center
  },
  % rounded callout / note (new)
  note/.style={
    rectangle, rounded corners=3pt,
    draw=clibord, line width=0.4pt, fill=clifill,
    inner sep=3pt, font=\sffamily\footnotesize, align=left
  },
  badge/.style={                                              % small inline marker
    rectangle, rounded corners=2pt,
    draw=hibord, line width=0.3pt, fill=hifill,
    inner sep=1.5pt, font=\sffamily\scriptsize\bfseries
  },
  % ─── edges (existing) ───
  inh/.style={                                  % UML inheritance
    -{Triangle[length=3mm,width=3mm,open]},
    draw=edgec, line width=0.5pt
  },
  agg/.style={                                  % UML aggregation
    {Diamond[length=2.5mm,width=2.5mm,open]}-{Stealth[length=2.2mm]},
    draw=edgec, line width=0.45pt
  },
  uses/.style={                                 % directed reference
    -{Stealth[length=2.2mm]}, draw=edgec, line width=0.5pt
  },
  arr/.style={                                  % plain arrow (graphs/trees)
    -{Stealth[length=2mm]}, draw=edgec, line width=0.45pt
  },
  edge/.style={                                 % undirected (trees)
    draw=edgec, line width=0.45pt
  },
  alt/.style={                                  % dashed: similar/alt
    -{Stealth[length=2mm]}, dashed, draw=edgesc, line width=0.45pt
  },
  ret/.style={                                  % dashed creates/returns
    -{Stealth[length=2mm]}, dashed, draw=edgesc, line width=0.45pt
  },
  thr/.style={                                  % threading link (red dashed)
    ->, dashed, draw=conbord, line width=0.5pt
  },
  % ─── new edges ───
  comp/.style={                                 % UML composition (filled diamond)
    {Diamond[length=2.5mm,width=2.5mm]}-{Stealth[length=2.2mm]},
    draw=edgec, line width=0.45pt
  },
  realize/.style={                              % UML realize (impl interface)
    -{Triangle[length=3mm,width=3mm,open]}, dashed,
    draw=edgec, line width=0.5pt
  },
  msgsync/.style={                              % synchronous message (sequence)
    -{Stealth[length=2.6mm]}, draw=edgec, line width=0.55pt
  },
  msgasync/.style={                             % asynchronous message
    -{Stealth[length=2.4mm,inset=0.4mm,fill=edgec]},
    draw=edgec, line width=0.5pt
  },
  msgret/.style={                               % return message (dashed)
    -{Stealth[length=2.2mm]}, dashed, draw=edgesc, line width=0.45pt
  },
  flow/.style={                                 % control flow (thicker, hi-vis)
    -{Stealth[length=2.4mm]}, draw=inkc, line width=0.7pt
  },
  focus/.style={                                % focus / hot path
    -{Stealth[length=2.4mm]}, draw=focusc, line width=0.8pt
  },
  selfloop/.style={                             % self loop on a node
    ->, draw=edgec, line width=0.45pt, in=120, out=60, loop, looseness=8
  },
  % ─── group / region (existing + new) ───
  group/.style={
    rectangle, rounded corners=4pt,
    draw=groupbord, line width=0.4pt,
    inner sep=10pt, fill=black!2
  },
  glabel/.style={font=\sffamily\bfseries\small, anchor=north west, inner sep=2pt},
  panel/.style={                                % stronger panel for layered diagrams
    rectangle, rounded corners=4pt,
    draw=neutbord, line width=0.45pt,
    inner sep=10pt, fill=neutfill!50
  },
  panelTitle/.style={font=\sffamily\bfseries\small, text=inkc, anchor=north west, inner sep=2pt},
  panelSub/.style={font=\sffamily\footnotesize, text=mutec, anchor=north west, inner sep=2pt},
  % ─── typography ───
  caption/.style={font=\sffamily\footnotesize\itshape, text=mutec, align=center},
  legendkey/.style={rectangle, minimum width=10mm, minimum height=3mm, inner sep=0pt, line width=0.3pt},
  % ─── grid / coordinate hints ───
  bggrid/.style={draw=gridc, line width=0.3pt, dotted},
}

% Korean font convenience macro — included by every Hangul diagram
\providecommand{\KOR}{}
     % adjust relative path
%   \begin{tikzpicture}[blog]  % apply 'blog' style
%   ...
%
% Backward-compat note: This file only ADDS to the public API. Existing styles
% (abs/con/imp/cli/hi/vx/q/inh/agg/uses/arr/edge/alt/ret/thr/group/glabel) are
% preserved. New helpers are documented below.

% ─── Core palette (fills paired with darker borders) ─────────
\definecolor{absfill}{RGB}{220,232,247}   % cool blue
\definecolor{absbord}{RGB}{72,118,182}
\definecolor{confill}{RGB}{249,222,222}   % warm red
\definecolor{conbord}{RGB}{197,93,93}
\definecolor{impfill}{RGB}{223,238,222}   % soft green
\definecolor{impbord}{RGB}{99,154,93}
\definecolor{clifill}{RGB}{251,239,205}   % pale yellow
\definecolor{clibord}{RGB}{200,167,72}
\definecolor{hifill} {RGB}{253,224,200}   % accent orange
\definecolor{hibord} {RGB}{210,135,68}
\definecolor{nodefill}{RGB}{250,250,250}
\definecolor{nodebord}{RGB}{120,120,120}
\definecolor{edgec}  {RGB}{80,80,80}
\definecolor{edgesc} {RGB}{145,145,145}
\definecolor{groupbord}{RGB}{200,200,200}

% ─── Extended palette (new) ──────────────────────────────────
\definecolor{prpfill}{RGB}{233,224,247}   % muted purple (state / interface)
\definecolor{prpbord}{RGB}{132,103,182}
\definecolor{tealfill}{RGB}{215,235,235}  % teal (services / utility)
\definecolor{tealbord}{RGB}{ 78,135,140}
\definecolor{warnfill}{RGB}{251,229,229}  % deeper warn (error states)
\definecolor{warnbord}{RGB}{180, 55, 55}
\definecolor{okfill}  {RGB}{220,240,225}  % success
\definecolor{okbord}  {RGB}{ 75,140, 90}
\definecolor{neutfill}{RGB}{238,238,238}  % neutral panel
\definecolor{neutbord}{RGB}{170,170,170}
\definecolor{inkc}    {RGB}{ 30, 30, 30}  % strong ink (titles)
\definecolor{mutec}   {RGB}{120,120,120}  % muted text
\definecolor{focusc}  {RGB}{210, 70, 70}  % focus / highlight accent
\definecolor{gridc}   {RGB}{220,220,220}  % background grid

\tikzset{
  blog/.style={
    font=\sffamily\small\linespread{1.18}\selectfont,
    every node/.style={font=\sffamily\small\linespread{1.18}\selectfont}
  },
  % ─── existing node roles (unchanged) ───
  cls/.style={
    rectangle, rounded corners=2pt,
    draw=nodebord, line width=0.4pt,
    minimum height=8mm, minimum width=28mm,
    inner sep=3pt, fill=nodefill, align=center
  },
  abs/.style={cls, draw=absbord, fill=absfill, font=\sffamily\itshape\small},
  con/.style={cls, draw=conbord, fill=confill},
  imp/.style={cls, draw=impbord, fill=impfill},
  cli/.style={cls, draw=clibord, fill=clifill},
  hi/.style={cls, draw=hibord, fill=hifill, font=\sffamily\bfseries\small},
  % ─── new node roles ───
  prp/.style={cls, draw=prpbord, fill=prpfill},               % protocol / interface / state
  tealn/.style={cls, draw=tealbord, fill=tealfill},           % service / utility
  warn/.style={cls, draw=warnbord, fill=warnfill, font=\sffamily\bfseries\small},  % error state
  ok/.style={cls, draw=okbord, fill=okfill},                  % success state
  neut/.style={cls, draw=neutbord, fill=neutfill},            % neutral
  % ─── size variants ───
  clsSm/.style={cls, minimum height=6mm, minimum width=20mm, font=\sffamily\footnotesize, inner sep=2pt},
  clsLg/.style={cls, minimum height=11mm, minimum width=36mm, font=\sffamily\normalsize},
  absSm/.style={clsSm, draw=absbord, fill=absfill, font=\sffamily\itshape\footnotesize},
  conSm/.style={clsSm, draw=conbord, fill=confill},
  impSm/.style={clsSm, draw=impbord, fill=impfill},
  % ─── circle (existing) ───
  vx/.style={
    circle, draw=absbord, line width=0.4pt,
    minimum size=8mm, fill=absfill, inner sep=0pt
  },
  vxc/.style={vx, draw=conbord, fill=confill},
  vxh/.style={vx, draw=hibord, fill=hifill, font=\sffamily\bfseries},
  % new circle variants
  vxi/.style={vx, draw=impbord, fill=impfill},                % impl green
  vxp/.style={vx, draw=prpbord, fill=prpfill},                % protocol purple
  vxw/.style={vx, draw=warnbord, fill=warnfill},              % warn red
  vxn/.style={vx, draw=nodebord, fill=nodefill},              % neutral
  vxSm/.style={vx, minimum size=6mm, font=\sffamily\footnotesize},
  vxLg/.style={vx, minimum size=11mm, font=\sffamily\normalsize},
  % ─── decision diamond (existing) ───
  q/.style={
    diamond, draw=clibord, line width=0.4pt, fill=clifill,
    inner sep=2pt, aspect=2.4, align=center
  },
  % rounded callout / note (new)
  note/.style={
    rectangle, rounded corners=3pt,
    draw=clibord, line width=0.4pt, fill=clifill,
    inner sep=3pt, font=\sffamily\footnotesize, align=left
  },
  badge/.style={                                              % small inline marker
    rectangle, rounded corners=2pt,
    draw=hibord, line width=0.3pt, fill=hifill,
    inner sep=1.5pt, font=\sffamily\scriptsize\bfseries
  },
  % ─── edges (existing) ───
  inh/.style={                                  % UML inheritance
    -{Triangle[length=3mm,width=3mm,open]},
    draw=edgec, line width=0.5pt
  },
  agg/.style={                                  % UML aggregation
    {Diamond[length=2.5mm,width=2.5mm,open]}-{Stealth[length=2.2mm]},
    draw=edgec, line width=0.45pt
  },
  uses/.style={                                 % directed reference
    -{Stealth[length=2.2mm]}, draw=edgec, line width=0.5pt
  },
  arr/.style={                                  % plain arrow (graphs/trees)
    -{Stealth[length=2mm]}, draw=edgec, line width=0.45pt
  },
  edge/.style={                                 % undirected (trees)
    draw=edgec, line width=0.45pt
  },
  alt/.style={                                  % dashed: similar/alt
    -{Stealth[length=2mm]}, dashed, draw=edgesc, line width=0.45pt
  },
  ret/.style={                                  % dashed creates/returns
    -{Stealth[length=2mm]}, dashed, draw=edgesc, line width=0.45pt
  },
  thr/.style={                                  % threading link (red dashed)
    ->, dashed, draw=conbord, line width=0.5pt
  },
  % ─── new edges ───
  comp/.style={                                 % UML composition (filled diamond)
    {Diamond[length=2.5mm,width=2.5mm]}-{Stealth[length=2.2mm]},
    draw=edgec, line width=0.45pt
  },
  realize/.style={                              % UML realize (impl interface)
    -{Triangle[length=3mm,width=3mm,open]}, dashed,
    draw=edgec, line width=0.5pt
  },
  msgsync/.style={                              % synchronous message (sequence)
    -{Stealth[length=2.6mm]}, draw=edgec, line width=0.55pt
  },
  msgasync/.style={                             % asynchronous message
    -{Stealth[length=2.4mm,inset=0.4mm,fill=edgec]},
    draw=edgec, line width=0.5pt
  },
  msgret/.style={                               % return message (dashed)
    -{Stealth[length=2.2mm]}, dashed, draw=edgesc, line width=0.45pt
  },
  flow/.style={                                 % control flow (thicker, hi-vis)
    -{Stealth[length=2.4mm]}, draw=inkc, line width=0.7pt
  },
  focus/.style={                                % focus / hot path
    -{Stealth[length=2.4mm]}, draw=focusc, line width=0.8pt
  },
  selfloop/.style={                             % self loop on a node
    ->, draw=edgec, line width=0.45pt, in=120, out=60, loop, looseness=8
  },
  % ─── group / region (existing + new) ───
  group/.style={
    rectangle, rounded corners=4pt,
    draw=groupbord, line width=0.4pt,
    inner sep=10pt, fill=black!2
  },
  glabel/.style={font=\sffamily\bfseries\small, anchor=north west, inner sep=2pt},
  panel/.style={                                % stronger panel for layered diagrams
    rectangle, rounded corners=4pt,
    draw=neutbord, line width=0.45pt,
    inner sep=10pt, fill=neutfill!50
  },
  panelTitle/.style={font=\sffamily\bfseries\small, text=inkc, anchor=north west, inner sep=2pt},
  panelSub/.style={font=\sffamily\footnotesize, text=mutec, anchor=north west, inner sep=2pt},
  % ─── typography ───
  caption/.style={font=\sffamily\footnotesize\itshape, text=mutec, align=center},
  legendkey/.style={rectangle, minimum width=10mm, minimum height=3mm, inner sep=0pt, line width=0.3pt},
  % ─── grid / coordinate hints ───
  bggrid/.style={draw=gridc, line width=0.3pt, dotted},
}

% Korean font convenience macro — included by every Hangul diagram
\providecommand{\KOR}{}

% _design-state.tex — TikZ state machine / FSM styles
%
% UML state diagram, protocol state machine, hardware FSM, kernel state graph
% 등에 공통. Mermaid의 state diagram은 *수동 정렬*이 보통 더 깔끔하므로
% TikZ로 직접 그리는 것이 표준 (CLAUDE.md 규칙).
%
% Companion to _design.tex — same palette and ink colors.
%
% Usage:
%   \documentclass[border=10pt,tikz]{standalone}
%   \usepackage{tikz}
%   \usetikzlibrary{positioning,arrows.meta,shapes,fit,calc,decorations.pathreplacing}
%   \begin{document}
%   % _design.tex — shared TikZ design system for all blog diagrams
% Inputs nothing on its own — meant to be \input{} from each diagram .tex
% AFTER \begin{document}, BEFORE \begin{tikzpicture}.
%
% Companion files:
%   _design-math.tex — math diagrams (PGFPlots, book-notes palette)
%
% Provides a unified visual language across GoF / DSA / EMC++ / Beautiful / 시리즈:
%   - Colors: muted, accessible, color-coded by role
%   - Node styles: consistent sizing, rounded, subtle borders
%   - Edge styles: UML-style inheritance, aggregation, plus dashed alternates
%   - Annotation: callouts, notes, badges
%
% Usage:
%   \documentclass[border=10pt,tikz]{standalone}
%   \usepackage{tikz}
%   \usetikzlibrary{positioning,arrows.meta,shapes,fit,backgrounds,calc,decorations.pathreplacing}
%   \begin{document}
%   % _design.tex — shared TikZ design system for all blog diagrams
% Inputs nothing on its own — meant to be \input{} from each diagram .tex
% AFTER \begin{document}, BEFORE \begin{tikzpicture}.
%
% Companion files:
%   _design-math.tex — math diagrams (PGFPlots, book-notes palette)
%
% Provides a unified visual language across GoF / DSA / EMC++ / Beautiful / 시리즈:
%   - Colors: muted, accessible, color-coded by role
%   - Node styles: consistent sizing, rounded, subtle borders
%   - Edge styles: UML-style inheritance, aggregation, plus dashed alternates
%   - Annotation: callouts, notes, badges
%
% Usage:
%   \documentclass[border=10pt,tikz]{standalone}
%   \usepackage{tikz}
%   \usetikzlibrary{positioning,arrows.meta,shapes,fit,backgrounds,calc,decorations.pathreplacing}
%   \begin{document}
%   \input{../_design.tex}     % adjust relative path
%   \begin{tikzpicture}[blog]  % apply 'blog' style
%   ...
%
% Backward-compat note: This file only ADDS to the public API. Existing styles
% (abs/con/imp/cli/hi/vx/q/inh/agg/uses/arr/edge/alt/ret/thr/group/glabel) are
% preserved. New helpers are documented below.

% ─── Core palette (fills paired with darker borders) ─────────
\definecolor{absfill}{RGB}{220,232,247}   % cool blue
\definecolor{absbord}{RGB}{72,118,182}
\definecolor{confill}{RGB}{249,222,222}   % warm red
\definecolor{conbord}{RGB}{197,93,93}
\definecolor{impfill}{RGB}{223,238,222}   % soft green
\definecolor{impbord}{RGB}{99,154,93}
\definecolor{clifill}{RGB}{251,239,205}   % pale yellow
\definecolor{clibord}{RGB}{200,167,72}
\definecolor{hifill} {RGB}{253,224,200}   % accent orange
\definecolor{hibord} {RGB}{210,135,68}
\definecolor{nodefill}{RGB}{250,250,250}
\definecolor{nodebord}{RGB}{120,120,120}
\definecolor{edgec}  {RGB}{80,80,80}
\definecolor{edgesc} {RGB}{145,145,145}
\definecolor{groupbord}{RGB}{200,200,200}

% ─── Extended palette (new) ──────────────────────────────────
\definecolor{prpfill}{RGB}{233,224,247}   % muted purple (state / interface)
\definecolor{prpbord}{RGB}{132,103,182}
\definecolor{tealfill}{RGB}{215,235,235}  % teal (services / utility)
\definecolor{tealbord}{RGB}{ 78,135,140}
\definecolor{warnfill}{RGB}{251,229,229}  % deeper warn (error states)
\definecolor{warnbord}{RGB}{180, 55, 55}
\definecolor{okfill}  {RGB}{220,240,225}  % success
\definecolor{okbord}  {RGB}{ 75,140, 90}
\definecolor{neutfill}{RGB}{238,238,238}  % neutral panel
\definecolor{neutbord}{RGB}{170,170,170}
\definecolor{inkc}    {RGB}{ 30, 30, 30}  % strong ink (titles)
\definecolor{mutec}   {RGB}{120,120,120}  % muted text
\definecolor{focusc}  {RGB}{210, 70, 70}  % focus / highlight accent
\definecolor{gridc}   {RGB}{220,220,220}  % background grid

\tikzset{
  blog/.style={
    font=\sffamily\small\linespread{1.18}\selectfont,
    every node/.style={font=\sffamily\small\linespread{1.18}\selectfont}
  },
  % ─── existing node roles (unchanged) ───
  cls/.style={
    rectangle, rounded corners=2pt,
    draw=nodebord, line width=0.4pt,
    minimum height=8mm, minimum width=28mm,
    inner sep=3pt, fill=nodefill, align=center
  },
  abs/.style={cls, draw=absbord, fill=absfill, font=\sffamily\itshape\small},
  con/.style={cls, draw=conbord, fill=confill},
  imp/.style={cls, draw=impbord, fill=impfill},
  cli/.style={cls, draw=clibord, fill=clifill},
  hi/.style={cls, draw=hibord, fill=hifill, font=\sffamily\bfseries\small},
  % ─── new node roles ───
  prp/.style={cls, draw=prpbord, fill=prpfill},               % protocol / interface / state
  tealn/.style={cls, draw=tealbord, fill=tealfill},           % service / utility
  warn/.style={cls, draw=warnbord, fill=warnfill, font=\sffamily\bfseries\small},  % error state
  ok/.style={cls, draw=okbord, fill=okfill},                  % success state
  neut/.style={cls, draw=neutbord, fill=neutfill},            % neutral
  % ─── size variants ───
  clsSm/.style={cls, minimum height=6mm, minimum width=20mm, font=\sffamily\footnotesize, inner sep=2pt},
  clsLg/.style={cls, minimum height=11mm, minimum width=36mm, font=\sffamily\normalsize},
  absSm/.style={clsSm, draw=absbord, fill=absfill, font=\sffamily\itshape\footnotesize},
  conSm/.style={clsSm, draw=conbord, fill=confill},
  impSm/.style={clsSm, draw=impbord, fill=impfill},
  % ─── circle (existing) ───
  vx/.style={
    circle, draw=absbord, line width=0.4pt,
    minimum size=8mm, fill=absfill, inner sep=0pt
  },
  vxc/.style={vx, draw=conbord, fill=confill},
  vxh/.style={vx, draw=hibord, fill=hifill, font=\sffamily\bfseries},
  % new circle variants
  vxi/.style={vx, draw=impbord, fill=impfill},                % impl green
  vxp/.style={vx, draw=prpbord, fill=prpfill},                % protocol purple
  vxw/.style={vx, draw=warnbord, fill=warnfill},              % warn red
  vxn/.style={vx, draw=nodebord, fill=nodefill},              % neutral
  vxSm/.style={vx, minimum size=6mm, font=\sffamily\footnotesize},
  vxLg/.style={vx, minimum size=11mm, font=\sffamily\normalsize},
  % ─── decision diamond (existing) ───
  q/.style={
    diamond, draw=clibord, line width=0.4pt, fill=clifill,
    inner sep=2pt, aspect=2.4, align=center
  },
  % rounded callout / note (new)
  note/.style={
    rectangle, rounded corners=3pt,
    draw=clibord, line width=0.4pt, fill=clifill,
    inner sep=3pt, font=\sffamily\footnotesize, align=left
  },
  badge/.style={                                              % small inline marker
    rectangle, rounded corners=2pt,
    draw=hibord, line width=0.3pt, fill=hifill,
    inner sep=1.5pt, font=\sffamily\scriptsize\bfseries
  },
  % ─── edges (existing) ───
  inh/.style={                                  % UML inheritance
    -{Triangle[length=3mm,width=3mm,open]},
    draw=edgec, line width=0.5pt
  },
  agg/.style={                                  % UML aggregation
    {Diamond[length=2.5mm,width=2.5mm,open]}-{Stealth[length=2.2mm]},
    draw=edgec, line width=0.45pt
  },
  uses/.style={                                 % directed reference
    -{Stealth[length=2.2mm]}, draw=edgec, line width=0.5pt
  },
  arr/.style={                                  % plain arrow (graphs/trees)
    -{Stealth[length=2mm]}, draw=edgec, line width=0.45pt
  },
  edge/.style={                                 % undirected (trees)
    draw=edgec, line width=0.45pt
  },
  alt/.style={                                  % dashed: similar/alt
    -{Stealth[length=2mm]}, dashed, draw=edgesc, line width=0.45pt
  },
  ret/.style={                                  % dashed creates/returns
    -{Stealth[length=2mm]}, dashed, draw=edgesc, line width=0.45pt
  },
  thr/.style={                                  % threading link (red dashed)
    ->, dashed, draw=conbord, line width=0.5pt
  },
  % ─── new edges ───
  comp/.style={                                 % UML composition (filled diamond)
    {Diamond[length=2.5mm,width=2.5mm]}-{Stealth[length=2.2mm]},
    draw=edgec, line width=0.45pt
  },
  realize/.style={                              % UML realize (impl interface)
    -{Triangle[length=3mm,width=3mm,open]}, dashed,
    draw=edgec, line width=0.5pt
  },
  msgsync/.style={                              % synchronous message (sequence)
    -{Stealth[length=2.6mm]}, draw=edgec, line width=0.55pt
  },
  msgasync/.style={                             % asynchronous message
    -{Stealth[length=2.4mm,inset=0.4mm,fill=edgec]},
    draw=edgec, line width=0.5pt
  },
  msgret/.style={                               % return message (dashed)
    -{Stealth[length=2.2mm]}, dashed, draw=edgesc, line width=0.45pt
  },
  flow/.style={                                 % control flow (thicker, hi-vis)
    -{Stealth[length=2.4mm]}, draw=inkc, line width=0.7pt
  },
  focus/.style={                                % focus / hot path
    -{Stealth[length=2.4mm]}, draw=focusc, line width=0.8pt
  },
  selfloop/.style={                             % self loop on a node
    ->, draw=edgec, line width=0.45pt, in=120, out=60, loop, looseness=8
  },
  % ─── group / region (existing + new) ───
  group/.style={
    rectangle, rounded corners=4pt,
    draw=groupbord, line width=0.4pt,
    inner sep=10pt, fill=black!2
  },
  glabel/.style={font=\sffamily\bfseries\small, anchor=north west, inner sep=2pt},
  panel/.style={                                % stronger panel for layered diagrams
    rectangle, rounded corners=4pt,
    draw=neutbord, line width=0.45pt,
    inner sep=10pt, fill=neutfill!50
  },
  panelTitle/.style={font=\sffamily\bfseries\small, text=inkc, anchor=north west, inner sep=2pt},
  panelSub/.style={font=\sffamily\footnotesize, text=mutec, anchor=north west, inner sep=2pt},
  % ─── typography ───
  caption/.style={font=\sffamily\footnotesize\itshape, text=mutec, align=center},
  legendkey/.style={rectangle, minimum width=10mm, minimum height=3mm, inner sep=0pt, line width=0.3pt},
  % ─── grid / coordinate hints ───
  bggrid/.style={draw=gridc, line width=0.3pt, dotted},
}

% Korean font convenience macro — included by every Hangul diagram
\providecommand{\KOR}{}
     % adjust relative path
%   \begin{tikzpicture}[blog]  % apply 'blog' style
%   ...
%
% Backward-compat note: This file only ADDS to the public API. Existing styles
% (abs/con/imp/cli/hi/vx/q/inh/agg/uses/arr/edge/alt/ret/thr/group/glabel) are
% preserved. New helpers are documented below.

% ─── Core palette (fills paired with darker borders) ─────────
\definecolor{absfill}{RGB}{220,232,247}   % cool blue
\definecolor{absbord}{RGB}{72,118,182}
\definecolor{confill}{RGB}{249,222,222}   % warm red
\definecolor{conbord}{RGB}{197,93,93}
\definecolor{impfill}{RGB}{223,238,222}   % soft green
\definecolor{impbord}{RGB}{99,154,93}
\definecolor{clifill}{RGB}{251,239,205}   % pale yellow
\definecolor{clibord}{RGB}{200,167,72}
\definecolor{hifill} {RGB}{253,224,200}   % accent orange
\definecolor{hibord} {RGB}{210,135,68}
\definecolor{nodefill}{RGB}{250,250,250}
\definecolor{nodebord}{RGB}{120,120,120}
\definecolor{edgec}  {RGB}{80,80,80}
\definecolor{edgesc} {RGB}{145,145,145}
\definecolor{groupbord}{RGB}{200,200,200}

% ─── Extended palette (new) ──────────────────────────────────
\definecolor{prpfill}{RGB}{233,224,247}   % muted purple (state / interface)
\definecolor{prpbord}{RGB}{132,103,182}
\definecolor{tealfill}{RGB}{215,235,235}  % teal (services / utility)
\definecolor{tealbord}{RGB}{ 78,135,140}
\definecolor{warnfill}{RGB}{251,229,229}  % deeper warn (error states)
\definecolor{warnbord}{RGB}{180, 55, 55}
\definecolor{okfill}  {RGB}{220,240,225}  % success
\definecolor{okbord}  {RGB}{ 75,140, 90}
\definecolor{neutfill}{RGB}{238,238,238}  % neutral panel
\definecolor{neutbord}{RGB}{170,170,170}
\definecolor{inkc}    {RGB}{ 30, 30, 30}  % strong ink (titles)
\definecolor{mutec}   {RGB}{120,120,120}  % muted text
\definecolor{focusc}  {RGB}{210, 70, 70}  % focus / highlight accent
\definecolor{gridc}   {RGB}{220,220,220}  % background grid

\tikzset{
  blog/.style={
    font=\sffamily\small\linespread{1.18}\selectfont,
    every node/.style={font=\sffamily\small\linespread{1.18}\selectfont}
  },
  % ─── existing node roles (unchanged) ───
  cls/.style={
    rectangle, rounded corners=2pt,
    draw=nodebord, line width=0.4pt,
    minimum height=8mm, minimum width=28mm,
    inner sep=3pt, fill=nodefill, align=center
  },
  abs/.style={cls, draw=absbord, fill=absfill, font=\sffamily\itshape\small},
  con/.style={cls, draw=conbord, fill=confill},
  imp/.style={cls, draw=impbord, fill=impfill},
  cli/.style={cls, draw=clibord, fill=clifill},
  hi/.style={cls, draw=hibord, fill=hifill, font=\sffamily\bfseries\small},
  % ─── new node roles ───
  prp/.style={cls, draw=prpbord, fill=prpfill},               % protocol / interface / state
  tealn/.style={cls, draw=tealbord, fill=tealfill},           % service / utility
  warn/.style={cls, draw=warnbord, fill=warnfill, font=\sffamily\bfseries\small},  % error state
  ok/.style={cls, draw=okbord, fill=okfill},                  % success state
  neut/.style={cls, draw=neutbord, fill=neutfill},            % neutral
  % ─── size variants ───
  clsSm/.style={cls, minimum height=6mm, minimum width=20mm, font=\sffamily\footnotesize, inner sep=2pt},
  clsLg/.style={cls, minimum height=11mm, minimum width=36mm, font=\sffamily\normalsize},
  absSm/.style={clsSm, draw=absbord, fill=absfill, font=\sffamily\itshape\footnotesize},
  conSm/.style={clsSm, draw=conbord, fill=confill},
  impSm/.style={clsSm, draw=impbord, fill=impfill},
  % ─── circle (existing) ───
  vx/.style={
    circle, draw=absbord, line width=0.4pt,
    minimum size=8mm, fill=absfill, inner sep=0pt
  },
  vxc/.style={vx, draw=conbord, fill=confill},
  vxh/.style={vx, draw=hibord, fill=hifill, font=\sffamily\bfseries},
  % new circle variants
  vxi/.style={vx, draw=impbord, fill=impfill},                % impl green
  vxp/.style={vx, draw=prpbord, fill=prpfill},                % protocol purple
  vxw/.style={vx, draw=warnbord, fill=warnfill},              % warn red
  vxn/.style={vx, draw=nodebord, fill=nodefill},              % neutral
  vxSm/.style={vx, minimum size=6mm, font=\sffamily\footnotesize},
  vxLg/.style={vx, minimum size=11mm, font=\sffamily\normalsize},
  % ─── decision diamond (existing) ───
  q/.style={
    diamond, draw=clibord, line width=0.4pt, fill=clifill,
    inner sep=2pt, aspect=2.4, align=center
  },
  % rounded callout / note (new)
  note/.style={
    rectangle, rounded corners=3pt,
    draw=clibord, line width=0.4pt, fill=clifill,
    inner sep=3pt, font=\sffamily\footnotesize, align=left
  },
  badge/.style={                                              % small inline marker
    rectangle, rounded corners=2pt,
    draw=hibord, line width=0.3pt, fill=hifill,
    inner sep=1.5pt, font=\sffamily\scriptsize\bfseries
  },
  % ─── edges (existing) ───
  inh/.style={                                  % UML inheritance
    -{Triangle[length=3mm,width=3mm,open]},
    draw=edgec, line width=0.5pt
  },
  agg/.style={                                  % UML aggregation
    {Diamond[length=2.5mm,width=2.5mm,open]}-{Stealth[length=2.2mm]},
    draw=edgec, line width=0.45pt
  },
  uses/.style={                                 % directed reference
    -{Stealth[length=2.2mm]}, draw=edgec, line width=0.5pt
  },
  arr/.style={                                  % plain arrow (graphs/trees)
    -{Stealth[length=2mm]}, draw=edgec, line width=0.45pt
  },
  edge/.style={                                 % undirected (trees)
    draw=edgec, line width=0.45pt
  },
  alt/.style={                                  % dashed: similar/alt
    -{Stealth[length=2mm]}, dashed, draw=edgesc, line width=0.45pt
  },
  ret/.style={                                  % dashed creates/returns
    -{Stealth[length=2mm]}, dashed, draw=edgesc, line width=0.45pt
  },
  thr/.style={                                  % threading link (red dashed)
    ->, dashed, draw=conbord, line width=0.5pt
  },
  % ─── new edges ───
  comp/.style={                                 % UML composition (filled diamond)
    {Diamond[length=2.5mm,width=2.5mm]}-{Stealth[length=2.2mm]},
    draw=edgec, line width=0.45pt
  },
  realize/.style={                              % UML realize (impl interface)
    -{Triangle[length=3mm,width=3mm,open]}, dashed,
    draw=edgec, line width=0.5pt
  },
  msgsync/.style={                              % synchronous message (sequence)
    -{Stealth[length=2.6mm]}, draw=edgec, line width=0.55pt
  },
  msgasync/.style={                             % asynchronous message
    -{Stealth[length=2.4mm,inset=0.4mm,fill=edgec]},
    draw=edgec, line width=0.5pt
  },
  msgret/.style={                               % return message (dashed)
    -{Stealth[length=2.2mm]}, dashed, draw=edgesc, line width=0.45pt
  },
  flow/.style={                                 % control flow (thicker, hi-vis)
    -{Stealth[length=2.4mm]}, draw=inkc, line width=0.7pt
  },
  focus/.style={                                % focus / hot path
    -{Stealth[length=2.4mm]}, draw=focusc, line width=0.8pt
  },
  selfloop/.style={                             % self loop on a node
    ->, draw=edgec, line width=0.45pt, in=120, out=60, loop, looseness=8
  },
  % ─── group / region (existing + new) ───
  group/.style={
    rectangle, rounded corners=4pt,
    draw=groupbord, line width=0.4pt,
    inner sep=10pt, fill=black!2
  },
  glabel/.style={font=\sffamily\bfseries\small, anchor=north west, inner sep=2pt},
  panel/.style={                                % stronger panel for layered diagrams
    rectangle, rounded corners=4pt,
    draw=neutbord, line width=0.45pt,
    inner sep=10pt, fill=neutfill!50
  },
  panelTitle/.style={font=\sffamily\bfseries\small, text=inkc, anchor=north west, inner sep=2pt},
  panelSub/.style={font=\sffamily\footnotesize, text=mutec, anchor=north west, inner sep=2pt},
  % ─── typography ───
  caption/.style={font=\sffamily\footnotesize\itshape, text=mutec, align=center},
  legendkey/.style={rectangle, minimum width=10mm, minimum height=3mm, inner sep=0pt, line width=0.3pt},
  % ─── grid / coordinate hints ───
  bggrid/.style={draw=gridc, line width=0.3pt, dotted},
}

% Korean font convenience macro — included by every Hangul diagram
\providecommand{\KOR}{}
        % core palette
%   % _design-state.tex — TikZ state machine / FSM styles
%
% UML state diagram, protocol state machine, hardware FSM, kernel state graph
% 등에 공통. Mermaid의 state diagram은 *수동 정렬*이 보통 더 깔끔하므로
% TikZ로 직접 그리는 것이 표준 (CLAUDE.md 규칙).
%
% Companion to _design.tex — same palette and ink colors.
%
% Usage:
%   \documentclass[border=10pt,tikz]{standalone}
%   \usepackage{tikz}
%   \usetikzlibrary{positioning,arrows.meta,shapes,fit,calc,decorations.pathreplacing}
%   \begin{document}
%   % _design.tex — shared TikZ design system for all blog diagrams
% Inputs nothing on its own — meant to be \input{} from each diagram .tex
% AFTER \begin{document}, BEFORE \begin{tikzpicture}.
%
% Companion files:
%   _design-math.tex — math diagrams (PGFPlots, book-notes palette)
%
% Provides a unified visual language across GoF / DSA / EMC++ / Beautiful / 시리즈:
%   - Colors: muted, accessible, color-coded by role
%   - Node styles: consistent sizing, rounded, subtle borders
%   - Edge styles: UML-style inheritance, aggregation, plus dashed alternates
%   - Annotation: callouts, notes, badges
%
% Usage:
%   \documentclass[border=10pt,tikz]{standalone}
%   \usepackage{tikz}
%   \usetikzlibrary{positioning,arrows.meta,shapes,fit,backgrounds,calc,decorations.pathreplacing}
%   \begin{document}
%   \input{../_design.tex}     % adjust relative path
%   \begin{tikzpicture}[blog]  % apply 'blog' style
%   ...
%
% Backward-compat note: This file only ADDS to the public API. Existing styles
% (abs/con/imp/cli/hi/vx/q/inh/agg/uses/arr/edge/alt/ret/thr/group/glabel) are
% preserved. New helpers are documented below.

% ─── Core palette (fills paired with darker borders) ─────────
\definecolor{absfill}{RGB}{220,232,247}   % cool blue
\definecolor{absbord}{RGB}{72,118,182}
\definecolor{confill}{RGB}{249,222,222}   % warm red
\definecolor{conbord}{RGB}{197,93,93}
\definecolor{impfill}{RGB}{223,238,222}   % soft green
\definecolor{impbord}{RGB}{99,154,93}
\definecolor{clifill}{RGB}{251,239,205}   % pale yellow
\definecolor{clibord}{RGB}{200,167,72}
\definecolor{hifill} {RGB}{253,224,200}   % accent orange
\definecolor{hibord} {RGB}{210,135,68}
\definecolor{nodefill}{RGB}{250,250,250}
\definecolor{nodebord}{RGB}{120,120,120}
\definecolor{edgec}  {RGB}{80,80,80}
\definecolor{edgesc} {RGB}{145,145,145}
\definecolor{groupbord}{RGB}{200,200,200}

% ─── Extended palette (new) ──────────────────────────────────
\definecolor{prpfill}{RGB}{233,224,247}   % muted purple (state / interface)
\definecolor{prpbord}{RGB}{132,103,182}
\definecolor{tealfill}{RGB}{215,235,235}  % teal (services / utility)
\definecolor{tealbord}{RGB}{ 78,135,140}
\definecolor{warnfill}{RGB}{251,229,229}  % deeper warn (error states)
\definecolor{warnbord}{RGB}{180, 55, 55}
\definecolor{okfill}  {RGB}{220,240,225}  % success
\definecolor{okbord}  {RGB}{ 75,140, 90}
\definecolor{neutfill}{RGB}{238,238,238}  % neutral panel
\definecolor{neutbord}{RGB}{170,170,170}
\definecolor{inkc}    {RGB}{ 30, 30, 30}  % strong ink (titles)
\definecolor{mutec}   {RGB}{120,120,120}  % muted text
\definecolor{focusc}  {RGB}{210, 70, 70}  % focus / highlight accent
\definecolor{gridc}   {RGB}{220,220,220}  % background grid

\tikzset{
  blog/.style={
    font=\sffamily\small\linespread{1.18}\selectfont,
    every node/.style={font=\sffamily\small\linespread{1.18}\selectfont}
  },
  % ─── existing node roles (unchanged) ───
  cls/.style={
    rectangle, rounded corners=2pt,
    draw=nodebord, line width=0.4pt,
    minimum height=8mm, minimum width=28mm,
    inner sep=3pt, fill=nodefill, align=center
  },
  abs/.style={cls, draw=absbord, fill=absfill, font=\sffamily\itshape\small},
  con/.style={cls, draw=conbord, fill=confill},
  imp/.style={cls, draw=impbord, fill=impfill},
  cli/.style={cls, draw=clibord, fill=clifill},
  hi/.style={cls, draw=hibord, fill=hifill, font=\sffamily\bfseries\small},
  % ─── new node roles ───
  prp/.style={cls, draw=prpbord, fill=prpfill},               % protocol / interface / state
  tealn/.style={cls, draw=tealbord, fill=tealfill},           % service / utility
  warn/.style={cls, draw=warnbord, fill=warnfill, font=\sffamily\bfseries\small},  % error state
  ok/.style={cls, draw=okbord, fill=okfill},                  % success state
  neut/.style={cls, draw=neutbord, fill=neutfill},            % neutral
  % ─── size variants ───
  clsSm/.style={cls, minimum height=6mm, minimum width=20mm, font=\sffamily\footnotesize, inner sep=2pt},
  clsLg/.style={cls, minimum height=11mm, minimum width=36mm, font=\sffamily\normalsize},
  absSm/.style={clsSm, draw=absbord, fill=absfill, font=\sffamily\itshape\footnotesize},
  conSm/.style={clsSm, draw=conbord, fill=confill},
  impSm/.style={clsSm, draw=impbord, fill=impfill},
  % ─── circle (existing) ───
  vx/.style={
    circle, draw=absbord, line width=0.4pt,
    minimum size=8mm, fill=absfill, inner sep=0pt
  },
  vxc/.style={vx, draw=conbord, fill=confill},
  vxh/.style={vx, draw=hibord, fill=hifill, font=\sffamily\bfseries},
  % new circle variants
  vxi/.style={vx, draw=impbord, fill=impfill},                % impl green
  vxp/.style={vx, draw=prpbord, fill=prpfill},                % protocol purple
  vxw/.style={vx, draw=warnbord, fill=warnfill},              % warn red
  vxn/.style={vx, draw=nodebord, fill=nodefill},              % neutral
  vxSm/.style={vx, minimum size=6mm, font=\sffamily\footnotesize},
  vxLg/.style={vx, minimum size=11mm, font=\sffamily\normalsize},
  % ─── decision diamond (existing) ───
  q/.style={
    diamond, draw=clibord, line width=0.4pt, fill=clifill,
    inner sep=2pt, aspect=2.4, align=center
  },
  % rounded callout / note (new)
  note/.style={
    rectangle, rounded corners=3pt,
    draw=clibord, line width=0.4pt, fill=clifill,
    inner sep=3pt, font=\sffamily\footnotesize, align=left
  },
  badge/.style={                                              % small inline marker
    rectangle, rounded corners=2pt,
    draw=hibord, line width=0.3pt, fill=hifill,
    inner sep=1.5pt, font=\sffamily\scriptsize\bfseries
  },
  % ─── edges (existing) ───
  inh/.style={                                  % UML inheritance
    -{Triangle[length=3mm,width=3mm,open]},
    draw=edgec, line width=0.5pt
  },
  agg/.style={                                  % UML aggregation
    {Diamond[length=2.5mm,width=2.5mm,open]}-{Stealth[length=2.2mm]},
    draw=edgec, line width=0.45pt
  },
  uses/.style={                                 % directed reference
    -{Stealth[length=2.2mm]}, draw=edgec, line width=0.5pt
  },
  arr/.style={                                  % plain arrow (graphs/trees)
    -{Stealth[length=2mm]}, draw=edgec, line width=0.45pt
  },
  edge/.style={                                 % undirected (trees)
    draw=edgec, line width=0.45pt
  },
  alt/.style={                                  % dashed: similar/alt
    -{Stealth[length=2mm]}, dashed, draw=edgesc, line width=0.45pt
  },
  ret/.style={                                  % dashed creates/returns
    -{Stealth[length=2mm]}, dashed, draw=edgesc, line width=0.45pt
  },
  thr/.style={                                  % threading link (red dashed)
    ->, dashed, draw=conbord, line width=0.5pt
  },
  % ─── new edges ───
  comp/.style={                                 % UML composition (filled diamond)
    {Diamond[length=2.5mm,width=2.5mm]}-{Stealth[length=2.2mm]},
    draw=edgec, line width=0.45pt
  },
  realize/.style={                              % UML realize (impl interface)
    -{Triangle[length=3mm,width=3mm,open]}, dashed,
    draw=edgec, line width=0.5pt
  },
  msgsync/.style={                              % synchronous message (sequence)
    -{Stealth[length=2.6mm]}, draw=edgec, line width=0.55pt
  },
  msgasync/.style={                             % asynchronous message
    -{Stealth[length=2.4mm,inset=0.4mm,fill=edgec]},
    draw=edgec, line width=0.5pt
  },
  msgret/.style={                               % return message (dashed)
    -{Stealth[length=2.2mm]}, dashed, draw=edgesc, line width=0.45pt
  },
  flow/.style={                                 % control flow (thicker, hi-vis)
    -{Stealth[length=2.4mm]}, draw=inkc, line width=0.7pt
  },
  focus/.style={                                % focus / hot path
    -{Stealth[length=2.4mm]}, draw=focusc, line width=0.8pt
  },
  selfloop/.style={                             % self loop on a node
    ->, draw=edgec, line width=0.45pt, in=120, out=60, loop, looseness=8
  },
  % ─── group / region (existing + new) ───
  group/.style={
    rectangle, rounded corners=4pt,
    draw=groupbord, line width=0.4pt,
    inner sep=10pt, fill=black!2
  },
  glabel/.style={font=\sffamily\bfseries\small, anchor=north west, inner sep=2pt},
  panel/.style={                                % stronger panel for layered diagrams
    rectangle, rounded corners=4pt,
    draw=neutbord, line width=0.45pt,
    inner sep=10pt, fill=neutfill!50
  },
  panelTitle/.style={font=\sffamily\bfseries\small, text=inkc, anchor=north west, inner sep=2pt},
  panelSub/.style={font=\sffamily\footnotesize, text=mutec, anchor=north west, inner sep=2pt},
  % ─── typography ───
  caption/.style={font=\sffamily\footnotesize\itshape, text=mutec, align=center},
  legendkey/.style={rectangle, minimum width=10mm, minimum height=3mm, inner sep=0pt, line width=0.3pt},
  % ─── grid / coordinate hints ───
  bggrid/.style={draw=gridc, line width=0.3pt, dotted},
}

% Korean font convenience macro — included by every Hangul diagram
\providecommand{\KOR}{}
        % core palette
%   % _design-state.tex — TikZ state machine / FSM styles
%
% UML state diagram, protocol state machine, hardware FSM, kernel state graph
% 등에 공통. Mermaid의 state diagram은 *수동 정렬*이 보통 더 깔끔하므로
% TikZ로 직접 그리는 것이 표준 (CLAUDE.md 규칙).
%
% Companion to _design.tex — same palette and ink colors.
%
% Usage:
%   \documentclass[border=10pt,tikz]{standalone}
%   \usepackage{tikz}
%   \usetikzlibrary{positioning,arrows.meta,shapes,fit,calc,decorations.pathreplacing}
%   \begin{document}
%   \input{../_design.tex}        % core palette
%   \input{../_design-state.tex}  % state-specific
%   \begin{tikzpicture}[blog, state]
%     \node[initial]            (i)   at (-2, 0)  {};
%     \node[stateBox]           (s1)  at ( 0, 0)  {Idle};
%     \node[stateBox]           (s2)  at ( 3, 0)  {Running};
%     \node[stateBox, accent]   (s3)  at ( 6, 0)  {Done};
%     \node[final]              (f)   at ( 8, 0)  {};
%
%     \draw[trans] (i)  -- (s1);
%     \draw[trans] (s1) -- node[transLabel]{start} (s2);
%     \draw[trans] (s2) -- node[transLabel]{finish} (s3);
%     \draw[trans] (s3) -- (f);
%
%     % self transition
%     \draw[trans] (s2) edge[loop above] node[transLabel]{tick} (s2);
%
%     % guard / action label
%     \draw[trans] (s1) to[bend right=30] node[transLabel]{stop\,/\,cleanup} (s3);
%
%     % choice (junction)
%     \node[choice] (c) at (4.5, -2) {};
%   \end{tikzpicture}

\tikzset{
  state/.style={
    every node/.append style={font=\sffamily\small},
  },
  % ─── State (rounded rectangle, default) ─────────────────
  stateBox/.style={
    rectangle, rounded corners=4pt,
    draw=nodebord, line width=0.5pt,
    minimum height=8mm, minimum width=22mm,
    inner sep=4pt, fill=nodefill, align=center,
  },
  stateBox/.append style={font=\sffamily\small},
  % Role-based variants — match _design.tex palette
  stateAbs/.style ={stateBox, draw=absbord, fill=absfill, font=\sffamily\itshape\small},
  stateCon/.style ={stateBox, draw=conbord, fill=confill},
  stateImp/.style ={stateBox, draw=impbord, fill=impfill},
  stateOk/.style  ={stateBox, draw=okbord,  fill=okfill},
  stateWarn/.style={stateBox, draw=warnbord, fill=warnfill, font=\sffamily\bfseries\small},
  statePrp/.style ={stateBox, draw=prpbord, fill=prpfill},   % protocol / interface
  stateNeut/.style={stateBox, draw=neutbord, fill=neutfill},
  accent/.style   ={draw=hibord, fill=hifill},               % apply on top (e.g., \node[stateBox, accent])
  % Compact / large variants
  stateSm/.style ={stateBox, minimum height=6mm, minimum width=16mm, font=\sffamily\footnotesize, inner sep=2pt},
  stateLg/.style ={stateBox, minimum height=11mm, minimum width=30mm},
  % ─── Initial / final / terminal ─────────────────────────
  initial/.style={
    circle, fill=inkc, draw=inkc, line width=0.5pt,
    minimum size=3.5mm, inner sep=0pt,
  },
  final/.style={                                              % bullseye (outer ring + inner dot)
    circle, draw=inkc, line width=0.5pt,
    minimum size=5mm, inner sep=0pt,
    path picture={
      \fill[inkc] (path picture bounding box.center) circle (1.2mm);
    },
  },
  terminate/.style={                                          % X (terminate pseudo-state)
    cross out, draw=warnbord, line width=0.8pt,
    minimum size=3mm, inner sep=0pt,
  },
  % ─── Choice / junction / fork-join ──────────────────────
  choice/.style={                                             % diamond decision
    diamond, draw=clibord, line width=0.45pt, fill=clifill,
    aspect=1, inner sep=1pt, minimum size=6mm,
    font=\sffamily\scriptsize, align=center,
  },
  junction/.style={                                           % small filled black circle (UML junction)
    circle, fill=inkc, draw=inkc, line width=0.4pt,
    minimum size=2.5mm, inner sep=0pt,
  },
  forkbar/.style={                                            % horizontal bar (fork / join)
    rectangle, fill=inkc, draw=inkc, line width=0.4pt,
    minimum width=20mm, minimum height=1.2mm, inner sep=0pt,
  },
  forkbarV/.style={forkbar, minimum width=1.2mm, minimum height=20mm},
  % ─── Transitions ────────────────────────────────────────
  trans/.style={
    -{Stealth[length=2.4mm]}, draw=edgec, line width=0.5pt,
  },
  transAlt/.style={                                           % dashed (lower priority / alt)
    -{Stealth[length=2.4mm]}, dashed, draw=edgesc, line width=0.5pt,
  },
  transFocus/.style={                                         % highlighted hot path
    -{Stealth[length=2.6mm]}, draw=focusc, line width=0.7pt,
  },
  transError/.style={                                         % error transition
    -{Stealth[length=2.4mm]}, draw=warnbord, line width=0.55pt,
  },
  % Label on transition — event[guard]/action 형식
  transLabel/.style={
    font=\sffamily\footnotesize, text=inkc,
    fill=white, fill opacity=0.85, text opacity=1,
    inner sep=2pt,
  },
  transLabelMute/.style={transLabel, text=mutec},
  % ─── Internal action listing (inside state box) ─────────
  internalAction/.style={
    font=\sffamily\footnotesize, text=inkc, anchor=north west, inner sep=2pt,
  },
  separator/.style={                                          % thin separator inside composite state
    draw=mutec, line width=0.3pt,
  },
  % ─── Composite / region (containing sub-states) ─────────
  composite/.style={
    rectangle, rounded corners=6pt,
    draw=nodebord, line width=0.5pt,
    inner sep=10pt, fill=neutfill!50,
  },
  compositeTitle/.style={
    font=\sffamily\bfseries\small, text=inkc,
    fill=neutfill, draw=neutbord, line width=0.3pt,
    inner sep=2pt, rounded corners=2pt,
    anchor=north west,
  },
  region/.style={                                             % parallel region inside composite
    rectangle, dashed,
    draw=mutec, line width=0.35pt,
    inner sep=8pt,
  },
}
  % state-specific
%   \begin{tikzpicture}[blog, state]
%     \node[initial]            (i)   at (-2, 0)  {};
%     \node[stateBox]           (s1)  at ( 0, 0)  {Idle};
%     \node[stateBox]           (s2)  at ( 3, 0)  {Running};
%     \node[stateBox, accent]   (s3)  at ( 6, 0)  {Done};
%     \node[final]              (f)   at ( 8, 0)  {};
%
%     \draw[trans] (i)  -- (s1);
%     \draw[trans] (s1) -- node[transLabel]{start} (s2);
%     \draw[trans] (s2) -- node[transLabel]{finish} (s3);
%     \draw[trans] (s3) -- (f);
%
%     % self transition
%     \draw[trans] (s2) edge[loop above] node[transLabel]{tick} (s2);
%
%     % guard / action label
%     \draw[trans] (s1) to[bend right=30] node[transLabel]{stop\,/\,cleanup} (s3);
%
%     % choice (junction)
%     \node[choice] (c) at (4.5, -2) {};
%   \end{tikzpicture}

\tikzset{
  state/.style={
    every node/.append style={font=\sffamily\small},
  },
  % ─── State (rounded rectangle, default) ─────────────────
  stateBox/.style={
    rectangle, rounded corners=4pt,
    draw=nodebord, line width=0.5pt,
    minimum height=8mm, minimum width=22mm,
    inner sep=4pt, fill=nodefill, align=center,
  },
  stateBox/.append style={font=\sffamily\small},
  % Role-based variants — match _design.tex palette
  stateAbs/.style ={stateBox, draw=absbord, fill=absfill, font=\sffamily\itshape\small},
  stateCon/.style ={stateBox, draw=conbord, fill=confill},
  stateImp/.style ={stateBox, draw=impbord, fill=impfill},
  stateOk/.style  ={stateBox, draw=okbord,  fill=okfill},
  stateWarn/.style={stateBox, draw=warnbord, fill=warnfill, font=\sffamily\bfseries\small},
  statePrp/.style ={stateBox, draw=prpbord, fill=prpfill},   % protocol / interface
  stateNeut/.style={stateBox, draw=neutbord, fill=neutfill},
  accent/.style   ={draw=hibord, fill=hifill},               % apply on top (e.g., \node[stateBox, accent])
  % Compact / large variants
  stateSm/.style ={stateBox, minimum height=6mm, minimum width=16mm, font=\sffamily\footnotesize, inner sep=2pt},
  stateLg/.style ={stateBox, minimum height=11mm, minimum width=30mm},
  % ─── Initial / final / terminal ─────────────────────────
  initial/.style={
    circle, fill=inkc, draw=inkc, line width=0.5pt,
    minimum size=3.5mm, inner sep=0pt,
  },
  final/.style={                                              % bullseye (outer ring + inner dot)
    circle, draw=inkc, line width=0.5pt,
    minimum size=5mm, inner sep=0pt,
    path picture={
      \fill[inkc] (path picture bounding box.center) circle (1.2mm);
    },
  },
  terminate/.style={                                          % X (terminate pseudo-state)
    cross out, draw=warnbord, line width=0.8pt,
    minimum size=3mm, inner sep=0pt,
  },
  % ─── Choice / junction / fork-join ──────────────────────
  choice/.style={                                             % diamond decision
    diamond, draw=clibord, line width=0.45pt, fill=clifill,
    aspect=1, inner sep=1pt, minimum size=6mm,
    font=\sffamily\scriptsize, align=center,
  },
  junction/.style={                                           % small filled black circle (UML junction)
    circle, fill=inkc, draw=inkc, line width=0.4pt,
    minimum size=2.5mm, inner sep=0pt,
  },
  forkbar/.style={                                            % horizontal bar (fork / join)
    rectangle, fill=inkc, draw=inkc, line width=0.4pt,
    minimum width=20mm, minimum height=1.2mm, inner sep=0pt,
  },
  forkbarV/.style={forkbar, minimum width=1.2mm, minimum height=20mm},
  % ─── Transitions ────────────────────────────────────────
  trans/.style={
    -{Stealth[length=2.4mm]}, draw=edgec, line width=0.5pt,
  },
  transAlt/.style={                                           % dashed (lower priority / alt)
    -{Stealth[length=2.4mm]}, dashed, draw=edgesc, line width=0.5pt,
  },
  transFocus/.style={                                         % highlighted hot path
    -{Stealth[length=2.6mm]}, draw=focusc, line width=0.7pt,
  },
  transError/.style={                                         % error transition
    -{Stealth[length=2.4mm]}, draw=warnbord, line width=0.55pt,
  },
  % Label on transition — event[guard]/action 형식
  transLabel/.style={
    font=\sffamily\footnotesize, text=inkc,
    fill=white, fill opacity=0.85, text opacity=1,
    inner sep=2pt,
  },
  transLabelMute/.style={transLabel, text=mutec},
  % ─── Internal action listing (inside state box) ─────────
  internalAction/.style={
    font=\sffamily\footnotesize, text=inkc, anchor=north west, inner sep=2pt,
  },
  separator/.style={                                          % thin separator inside composite state
    draw=mutec, line width=0.3pt,
  },
  % ─── Composite / region (containing sub-states) ─────────
  composite/.style={
    rectangle, rounded corners=6pt,
    draw=nodebord, line width=0.5pt,
    inner sep=10pt, fill=neutfill!50,
  },
  compositeTitle/.style={
    font=\sffamily\bfseries\small, text=inkc,
    fill=neutfill, draw=neutbord, line width=0.3pt,
    inner sep=2pt, rounded corners=2pt,
    anchor=north west,
  },
  region/.style={                                             % parallel region inside composite
    rectangle, dashed,
    draw=mutec, line width=0.35pt,
    inner sep=8pt,
  },
}
  % state-specific
%   \begin{tikzpicture}[blog, state]
%     \node[initial]            (i)   at (-2, 0)  {};
%     \node[stateBox]           (s1)  at ( 0, 0)  {Idle};
%     \node[stateBox]           (s2)  at ( 3, 0)  {Running};
%     \node[stateBox, accent]   (s3)  at ( 6, 0)  {Done};
%     \node[final]              (f)   at ( 8, 0)  {};
%
%     \draw[trans] (i)  -- (s1);
%     \draw[trans] (s1) -- node[transLabel]{start} (s2);
%     \draw[trans] (s2) -- node[transLabel]{finish} (s3);
%     \draw[trans] (s3) -- (f);
%
%     % self transition
%     \draw[trans] (s2) edge[loop above] node[transLabel]{tick} (s2);
%
%     % guard / action label
%     \draw[trans] (s1) to[bend right=30] node[transLabel]{stop\,/\,cleanup} (s3);
%
%     % choice (junction)
%     \node[choice] (c) at (4.5, -2) {};
%   \end{tikzpicture}

\tikzset{
  state/.style={
    every node/.append style={font=\sffamily\small},
  },
  % ─── State (rounded rectangle, default) ─────────────────
  stateBox/.style={
    rectangle, rounded corners=4pt,
    draw=nodebord, line width=0.5pt,
    minimum height=8mm, minimum width=22mm,
    inner sep=4pt, fill=nodefill, align=center,
  },
  stateBox/.append style={font=\sffamily\small},
  % Role-based variants — match _design.tex palette
  stateAbs/.style ={stateBox, draw=absbord, fill=absfill, font=\sffamily\itshape\small},
  stateCon/.style ={stateBox, draw=conbord, fill=confill},
  stateImp/.style ={stateBox, draw=impbord, fill=impfill},
  stateOk/.style  ={stateBox, draw=okbord,  fill=okfill},
  stateWarn/.style={stateBox, draw=warnbord, fill=warnfill, font=\sffamily\bfseries\small},
  statePrp/.style ={stateBox, draw=prpbord, fill=prpfill},   % protocol / interface
  stateNeut/.style={stateBox, draw=neutbord, fill=neutfill},
  accent/.style   ={draw=hibord, fill=hifill},               % apply on top (e.g., \node[stateBox, accent])
  % Compact / large variants
  stateSm/.style ={stateBox, minimum height=6mm, minimum width=16mm, font=\sffamily\footnotesize, inner sep=2pt},
  stateLg/.style ={stateBox, minimum height=11mm, minimum width=30mm},
  % ─── Initial / final / terminal ─────────────────────────
  initial/.style={
    circle, fill=inkc, draw=inkc, line width=0.5pt,
    minimum size=3.5mm, inner sep=0pt,
  },
  final/.style={                                              % bullseye (outer ring + inner dot)
    circle, draw=inkc, line width=0.5pt,
    minimum size=5mm, inner sep=0pt,
    path picture={
      \fill[inkc] (path picture bounding box.center) circle (1.2mm);
    },
  },
  terminate/.style={                                          % X (terminate pseudo-state)
    cross out, draw=warnbord, line width=0.8pt,
    minimum size=3mm, inner sep=0pt,
  },
  % ─── Choice / junction / fork-join ──────────────────────
  choice/.style={                                             % diamond decision
    diamond, draw=clibord, line width=0.45pt, fill=clifill,
    aspect=1, inner sep=1pt, minimum size=6mm,
    font=\sffamily\scriptsize, align=center,
  },
  junction/.style={                                           % small filled black circle (UML junction)
    circle, fill=inkc, draw=inkc, line width=0.4pt,
    minimum size=2.5mm, inner sep=0pt,
  },
  forkbar/.style={                                            % horizontal bar (fork / join)
    rectangle, fill=inkc, draw=inkc, line width=0.4pt,
    minimum width=20mm, minimum height=1.2mm, inner sep=0pt,
  },
  forkbarV/.style={forkbar, minimum width=1.2mm, minimum height=20mm},
  % ─── Transitions ────────────────────────────────────────
  trans/.style={
    -{Stealth[length=2.4mm]}, draw=edgec, line width=0.5pt,
  },
  transAlt/.style={                                           % dashed (lower priority / alt)
    -{Stealth[length=2.4mm]}, dashed, draw=edgesc, line width=0.5pt,
  },
  transFocus/.style={                                         % highlighted hot path
    -{Stealth[length=2.6mm]}, draw=focusc, line width=0.7pt,
  },
  transError/.style={                                         % error transition
    -{Stealth[length=2.4mm]}, draw=warnbord, line width=0.55pt,
  },
  % Label on transition — event[guard]/action 형식
  transLabel/.style={
    font=\sffamily\footnotesize, text=inkc,
    fill=white, fill opacity=0.85, text opacity=1,
    inner sep=2pt,
  },
  transLabelMute/.style={transLabel, text=mutec},
  % ─── Internal action listing (inside state box) ─────────
  internalAction/.style={
    font=\sffamily\footnotesize, text=inkc, anchor=north west, inner sep=2pt,
  },
  separator/.style={                                          % thin separator inside composite state
    draw=mutec, line width=0.3pt,
  },
  % ─── Composite / region (containing sub-states) ─────────
  composite/.style={
    rectangle, rounded corners=6pt,
    draw=nodebord, line width=0.5pt,
    inner sep=10pt, fill=neutfill!50,
  },
  compositeTitle/.style={
    font=\sffamily\bfseries\small, text=inkc,
    fill=neutfill, draw=neutbord, line width=0.3pt,
    inner sep=2pt, rounded corners=2pt,
    anchor=north west,
  },
  region/.style={                                             % parallel region inside composite
    rectangle, dashed,
    draw=mutec, line width=0.35pt,
    inner sep=8pt,
  },
}

\begin{tikzpicture}[blog, state, x=1cm, y=1cm]

  % ───────────────── future panel (위) ─────────────────
  \node[panelTitle] at (-0.2, 5.4) {\texttt{future} — 백그라운드 계산의 약속};

  \node[initial]                (fs)    at (0.0, 4.0) {};
  \node[stateBox]               (frun)  at (3.0, 4.0) {Running};
  \node[stateBox, draw=okbord, fill=okfill] (fdone) at (8.0, 4.0)
        {Realized\\[2pt]\textit{값 결정됨}};
  \node[stateWarn]              (fcanc) at (5.5, 2.0) {Cancelled};

  \draw[trans] (fs)   -- node[transLabel]{\texttt{(future expr)}} (frun);
  \draw[trans] (frun) -- node[transLabel]{계산 완료} (fdone);
  \draw[transAlt] (frun.south) to[bend right=15]
        node[transLabel, near end]{\texttt{future-cancel}} (fcanc.west);

  \draw[trans] (fdone.north) to[loop above]
        node[transLabel]{\texttt{@f} = 같은 값 (immutable)} (fdone.north);

  % ───────────────── promise panel (아래) ─────────────────
  \node[panelTitle] at (-0.2, 0.4) {\texttt{promise} — 외부가 채워 넣을 약속};

  \node[initial]                  (ps)    at (0.0, -0.8) {};
  \node[stateBox]                 (ppend) at (3.0, -0.8) {Pending\\[2pt]\textit{값 없음}};
  \node[stateBox, draw=okbord, fill=okfill] (pdel) at (8.0, -0.8)
        {Delivered\\[2pt]\textit{값 결정됨}};

  \draw[trans] (ps)    -- node[transLabel]{\texttt{(promise)}} (ppend);
  \draw[trans] (ppend) -- node[transLabel]{\texttt{(deliver p v)}\\첫 호출만 적용} (pdel);

  % 무시되는 두 번째 deliver
  \draw[transError] (pdel.north) to[loop above]
        node[transLabel, text=warnbord]{두 번째 \texttt{deliver} — 무시} (pdel.north);

  % deref always blocks until delivered
  \draw[transAlt] (ppend.south) to[bend right=30]
        node[transLabel, below]{\texttt{@p} 호출자 블로킹}
        (pdel.south);

  % ───────────────── 비교 노트 ─────────────────
  \node[note, text width=140mm, align=left, anchor=north]
       at (5.0, -3.6)
       {\textbf{핵심 차이}.\\[2pt]
        \texttt{future}: 생성 즉시 계산 시작 — 결과는 비동기로 \textit{도착}한다.\\[2pt]
        \texttt{promise}: 누군가 \texttt{deliver}하기 전까지는 값이 없음 — 스레드 간 \textit{시그널} 전용.\\[2pt]
        두 경우 모두 \textbf{한 번 결정된 값은 immutable} — \texttt{@}로 몇 번을 받아도 같은 snapshot.};

\end{tikzpicture}
\end{document}
